\documentclass[12pt,a4paper]{article}
\usepackage{fontspec}
\setmainfont{DejaVu Serif}
\setsansfont{DejaVu Sans}
\setmonofont{DejaVu Sans Mono}[Scale=0.85]
\usepackage[margin=2.2cm]{geometry}
\usepackage{amsmath}
\usepackage{amssymb}
\usepackage{xcolor}
\usepackage{enumitem}
\usepackage{booktabs}
\usepackage{array}
\usepackage{listings}
\usepackage{hyperref}
\usepackage{titlesec}

\titleformat{\section}{\large\bfseries}{}{0pt}{}
\setlength{\parskip}{0.5em}
\setlength{\parindent}{0pt}
\sloppy
\setlength{\emergencystretch}{3em}
\hypersetup{hidelinks}

\lstset{
  basicstyle=\ttfamily\small,
  keywordstyle=\color{blue!70!black}\bfseries,
  commentstyle=\color{green!45!black},
  frame=single, framerule=0.3pt, rulecolor=\color{gray!60},
  backgroundcolor=\color{gray!6},
  columns=fullflexible, keepspaces=true,
  language=C, showstringspaces=false,
  xleftmargin=1em, xrightmargin=1em
}

\newcommand{\question}[2]{\section{Дасгал #1}\textbf{Асуулт.} #2\par\medskip\textbf{Хариулт.}\par}
\newcommand{\vect}[1]{(#1)}

\begin{document}

\begin{center}
  {\Large\bfseries Лабораторийн ажил №7}\\[0.4em]
  {\large Түгжрэл (Deadlocks) — Бүлэг 8, Practice Exercises}\\[0.8em]
  F.CSM302 Үйлдлийн систем\\
  Гүйцэтгэсэн: Б.Билгүүн (B190910014)\\[0.3em]
  Сурах бичиг: A. Silberschatz, P. B. Galvin, G. Gagne, \textit{Operating System Concepts}, 10th ed., Wiley, 2018
\end{center}
\hrule
\bigskip

Түгжрэл үүсэх зайлшгүй дөрвөн нөхцөл (§8.3.1): \textbf{mutual exclusion} (нөөцийг нэг удаад нэг урсгал ашиглана), \textbf{hold and wait} (нөөц эзэмшсээр өөр нөөц хүлээнэ), \textbf{no preemption} (нөөцийг хүчээр булааж авахгүй), \textbf{circular wait} (тойрог хүлээлт). Дөрвүүлээ зэрэг биелбэл түгжрэл үүсэх боломжтой.

%---------------------------------------------------------------
\question{8.1}{Компьютерийн системтэй холбоогүй түгжрэлийн гурван жишээ дурд.}

\begin{enumerate}[leftmargin=1.5em]
  \item \textbf{Уулзвар дахь машинууд.} Гэрлэн дохиогүй дөрвөн замын уулзварт дөрвөн машин зэрэг орж ирээд тус бүр баруун талынхаа машиныг өнгөрөхийг хүлээвэл хэн нь ч хөдлөхгүй (тойрог хүлээлт). Зам (уулзварын хэсэг) нь дундаа хэрэглэх боломжгүй нөөц.
  \item \textbf{Нэг эгнээтэй гүүр.} Эсрэг талаас хоёр машин гүүрийн дунд нүүр тулбал аль аль нь ухрахгүй бол хоёулаа зогсоно (Exercise 8.30-ын санаа).
  \item \textbf{Хоёр хүн, хоёр хэрэгсэл.} Нэг хүн сав, нөгөө нь халбага барьж аль аль нь нөгөөгийнхөө барьсныг хүлээвэл хоол хийгдэхгүй. Мөн хоёр хүн утсаар нэгнээ залгахад хоёулаа "завгүй" дуут дохио сонсож, нөгөө нь тасалж залгахыг хүлээх.
\end{enumerate}

%---------------------------------------------------------------
\question{8.2}{Систем аюулгүй биш төлөвт байна гэж үзье. Урсгалууд түгжрэлд орохгүйгээр гүйцэтгэлээ дуусгах боломжтойг харуул.}

Аюулгүй биш төлөв гэдэг нь \emph{хамгийн муу тохиолдолд} (урсгал бүр \textit{Max}-аа нэг дор шаардвал) түгжрэл үүсч болзошгүй гэсэн үг; түгжрэл заавал үүснэ гэсэн үг биш. \textit{Max} нь зөвхөн дээд хязгаар бөгөөд урсгал түүнээс бага нөөц шаардаад дуусч болно, эсвэл нөөцөө эрт чөлөөлж болно.

\textbf{Жишээ.} 12 соронзон хальстай систем (§8.6.2): $T_0$ — Max 10, эзэмшиж буй 5; $T_1$ — Max 4, эзэмшиж буй 2; $T_2$ — Max 9, эзэмшиж буй 3; чөлөөтэй 2. Энэ төлөв аюулгүй биш: $T_1$ дуусаад 4 чөлөөтэй болсон ч $T_0$ (5 хэрэгтэй) болон $T_2$ (6 хэрэгтэй)-ын аль нь ч Max хүртлээ шаардвал хүрэлцэхгүй. Гэвч $T_2$ нэмэлт нөөц шаардалгүй ажлаа дуусгаж 3 хальсаа чөлөөлбөл (Max нь зөвхөн дээд хязгаар) чөлөөтэй 7 болж, $T_0$ 5 хальс авч дуусна. Ингээд бүх урсгал түгжрэлгүй дуусна. Иймд аюулгүй биш төлөвөөс түгжрэлгүй дуусах гүйцэтгэлийн дараалал байж болно.

%---------------------------------------------------------------
\question{8.3}{Системийн агшин зураг (A B C D нөөц):
\begin{center}
\begin{tabular}{lccc}
\toprule
 & Allocation & Max & Available\\
\midrule
$T_0$ & 0 0 1 2 & 0 0 1 2 & 1 5 2 0\\
$T_1$ & 1 0 0 0 & 1 7 5 0 & \\
$T_2$ & 1 3 5 4 & 2 3 5 6 & \\
$T_3$ & 0 6 3 2 & 0 6 5 2 & \\
$T_4$ & 0 0 1 4 & 0 6 5 6 & \\
\bottomrule
\end{tabular}
\end{center}
Банкирын алгоритмаар: (a) Need матриц; (b) систем аюулгүй төлөвт байгаа эсэх; (c) $T_1$-ээс $(0,4,2,0)$ хүсэлт ирвэл шууд хангаж болох уу?}

\textbf{(a)} $\text{Need} = \text{Max} - \text{Allocation}$:
\begin{center}
\begin{tabular}{lc}
\toprule
 & Need (A B C D)\\
\midrule
$T_0$ & 0 0 0 0\\
$T_1$ & 0 7 5 0\\
$T_2$ & 1 0 0 2\\
$T_3$ & 0 0 2 0\\
$T_4$ & 0 6 4 2\\
\bottomrule
\end{tabular}
\end{center}

\textbf{(b)} Аюулгүй байдлын алгоритм, $\text{Work} = \text{Available} = (1,5,2,0)$:
\begin{center}
\begin{tabular}{lll}
\toprule
Алхам & Нөхцөл $\text{Need}_i \le \text{Work}$ & Шинэ Work $=$ Work $+$ Allocation$_i$\\
\midrule
$T_0$ & $(0,0,0,0) \le (1,5,2,0)$ & $(1,5,3,2)$\\
$T_2$ & $(1,0,0,2) \le (1,5,3,2)$ & $(2,8,8,6)$\\
$T_3$ & $(0,0,2,0) \le (2,8,8,6)$ & $(2,14,11,8)$\\
$T_4$ & $(0,6,4,2) \le (2,14,11,8)$ & $(2,14,12,12)$\\
$T_1$ & $(0,7,5,0) \le (2,14,12,12)$ & $(3,14,12,12)$\\
\bottomrule
\end{tabular}
\end{center}
Бүх урсгал дуусах $\langle T_0, T_2, T_3, T_4, T_1\rangle$ дараалал байгаа тул систем \textbf{аюулгүйн төлөвт} байна.

\textbf{(c)} Хүсэлтийн алгоритм (§8.6.3.2): $\text{Request}_1 = (0,4,2,0) \le \text{Need}_1 = (0,7,5,0)$ $\checkmark$ ба $\le \text{Available} = (1,5,2,0)$ $\checkmark$. Хүсэлтийг өгсөн гэж үзвэл: Available $= (1,1,0,0)$, Allocation$_1 = (1,4,2,0)$, Need$_1 = (0,3,3,0)$. Аюулгүй байдлын алгоритм:
$T_0 \to (1,1,1,2)$, $T_2 \to (2,4,6,6)$, $T_3 \to (2,10,9,8)$, $T_4 \to (2,10,10,12)$, $T_1 \to (3,14,12,12)$.
Шинэ төлөв аюулгүй ($\langle T_0, T_2, T_3, T_4, T_1\rangle$) тул хүсэлтийг \textbf{шууд хангаж болно}. (Тооцоог \texttt{lab7/banker.py}-ээр шалгасан.)

%---------------------------------------------------------------
\question{8.4}{Түгжрэлээс сэргийлэх нэг арга: бусад бүх нөөцийн өмнө авах ёстой нэг "дээд эрэмбийн" нөөц F нэмэх (containment). A…E объектын аль нэгийн түгжээг авахын өмнө F-ийн түгжээг заавал авна. Энэ схемийг §8.5.4-ийн circular-wait схемтэй харьцуул.}

\textbf{Containment (F түгжээ).} Нэг урсгал F-ийг барьж байх үед өөр ямар ч урсгал A…E-ийн түгжээг авч чадахгүй тул A…E-ийн хувьд тойрог хүлээлт үүсэх боломжгүй — түгжрэлгүй. Гэвч F нь бүх зүйлийг \emph{дараалалд} оруулдаг "том түгжээ" (coarse-grained lock) болно: A-тай ажиллах урсгал ба D-тэй ажиллах урсгал хамааралгүй байсан ч зэрэг ажиллаж чадахгүй. Зэрэгцээ ажиллагаа эрс буурч, F-ийн төлөөх өрсөлдөөн (contention) гацаа үүсгэнэ. Мөн урсгал хэр удаан F-ийг барихаас хамаарч хүлээлт урт болно.

\textbf{Circular-wait схем (§8.5.4).} Нөөц бүрт дугаар оноож, урсгалууд түгжээг зөвхөн өсөх дарааллаар авна. Түгжрэлгүй байдлыг мөн баталгаажуулна, мөн хамааралгүй нөөцүүдийг зэрэг ашиглах боломжтой тул \textbf{зэрэгцээ ажиллагаа илүү}. Сул тал нь: бүх кодыг дарааллын дүрмийг мөрдүүлэх хэцүү (жнь Figure 8.7-ийн \texttt{transaction()} шиг түгжээ динамикаар сонгогдох үед), том системд дарааллыг тогтоох, хадгалах ажил их, Linux-ийн \texttt{lockdep} шиг хэрэгслээр шалгах шаардлагатай.

\textbf{Дүгнэлт:} containment нь энгийн, алдаа гаргахад хэцүү боловч зэрэгцээ байдлыг алддаг; эрэмбэлсэн түгжээ нь гүйцэтгэл сайн боловч хэрэгжүүлэх, засварлахад сахилга бат шаарддаг.

%---------------------------------------------------------------
\question{8.5}{§8.6.3-ийн аюулгүй байдлын алгоритм $m \times n^2$ эрэмбийн үйлдэл шаарддагийг батал.}

$n$ — урсгалын тоо, $m$ — нөөцийн төрлийн тоо.
\begin{enumerate}[leftmargin=1.5em]
  \item Эхлүүлэх: $\text{Work} = \text{Available}$ ($m$ үйлдэл), $\text{Finish}[i] = \text{false}$ ($n$ үйлдэл) — $O(m+n)$.
  \item Нэг давталт: \textit{Finish}$[i]=$false ба $\text{Need}_i \le \text{Work}$ байх $i$-г хайна. Урсгал бүрийн хувьд хоёр вектор харьцуулах нь $m$ үйлдэл, $n$ урсгал шалгана → $O(m \cdot n)$. Олдсон урсгалын хувьд $\text{Work} \mathrel{+}= \text{Allocation}_i$ — $O(m)$.
  \item Давталт бүрт дор хаяж нэг урсгал дуусна (эс бөгөөс алгоритм зогсоно), тиймээс давталт хамгийн ихдээ $n$ удаа хийгдэнэ.
\end{enumerate}
Нийт: $n \cdot (m \cdot n + m) = O(m \cdot n^2)$. Хамгийн муу тохиолдол: давталт бүрт олдох урсгал жагсаалтын хамгийн сүүлд байх үед бүх $n$ урсгалыг дахин шалгах шаардлагатай тул энэ дээд хязгаарт хүрнэ.

%---------------------------------------------------------------
\question{8.6}{Сард 5000 ажил гүйцэтгэдэг, түгжрэлээс сэргийлэх/гарах схемгүй систем. Сард 2 удаа түгжрэл үүсч, удаа бүр ~10 ажлыг зогсоож дахин ажиллуулдаг. Ажил бүр ~\$2 үнэтэй, зогсоосон ажлууд ойролцоогоор хагас гүйцэтгэгдсэн байдаг. Банкирын шиг түгжрэлээс гарах алгоритм суулгавал ажлын дундаж гүйцэтгэлийн хугацаа ~10\%-иар өснө; машин 30\% сул тул бүх ажил гүйцэтгэгдэх ч эргэлтийн хугацаа дунджаар ~20\%-иар өснө. (a) Суулгахын талд, (b) эсрэг талд ямар үндэслэл байна?}

\textbf{Алдагдлын тооцоо:} сард $2 \times 10 = 20$ ажил дахин ажиллана, хагас гүйцэтгэгдсэн тул алдагдал $20 \times \$2 \times 0.5 = \$20$/сар. Харин алгоритмын нэмэлт ачаалал: $5000 \times \$2 \times 10\% = \$1000$/сарын CPU цаг.

\textbf{(a) Суулгахын талд:}
\begin{itemize}[leftmargin=1.5em]
  \item Түгжрэл огт үүсэхгүй — оператор гараар илрүүлж, ажил зогсоож, дахин ажиллуулах шаардлагагүй (хүний цаг, алдаа гарах эрсдэл буурна).
  \item Урьдчилан таамаглах боломжтой байдал: зогсоосон ажлын хагас дутуу үр дүн (файл, өгөгдлийн сангийн өөрчлөлт)-ийг буцаах асуудал үүсэхгүй.
  \item Машин 30\% сул тул нэмэлт 10\% ачааллыг даана, бүх ажил гүйцэтгэгдсээр байна.
\end{itemize}

\textbf{(b) Эсрэг талд:}
\begin{itemize}[leftmargin=1.5em]
  \item Түгжрэлийн алдагдал ердөө $\approx\$20$/сар, харин алгоритмын зардал $\approx\$1000$/сар — эдийн засгийн хувьд ашиггүй.
  \item Эргэлтийн хугацаа бүх 5000 ажилд 20\% өснө — сард 20 ажлын саатлаас хамаагүй их хэрэглэгчид нөлөөлнө.
  \item Банкирын алгоритм урсгал бүр \textit{Max}-аа урьдчилан мэдэгдэхийг шаарддаг; бодит ажлуудад үүнийг тодорхойлоход хэцүү, нөөцийг хэтрүүлэн нөөцөлж ашиглалтыг бууруулна.
  \item Сул 30\% нь ирээдүйн өсөлтөд зориулсан нөөц байж болох тул түүнийг алгоритмд зарцуулах нь эрсдэлтэй.
\end{itemize}
Энэ тохиолдолд түгжрэл ховор, алдагдал бага тул ихэнх ерөнхий зориулалтын OS (Linux, Windows)-ийн адил \emph{түгжрэлийг үл тоох / илрүүлж сэргээх} хандлага илүү тохиромжтой.

%---------------------------------------------------------------
\question{8.7}{Систем зарим урсгалаа хомсдолд (starvation) орсныг илрүүлж чадах уу? "Тийм" бол хэрхэн, "үгүй" бол хомсдлыг хэрхэн шийдвэрлэх вэ?}

\textbf{Ерөнхийдөө үгүй — найдвартай илрүүлэх боломжгүй.} Хомсдол гэдэг нь урсгал \emph{хязгааргүй} удаан хүлээх явдал; тодорхой агшинд урсгал удаан хүлээж байгааг харж болох ч удахгүй нөөц авах эсэхийг мэдэх нь ирээдүйн хуваарилалтыг таамаглахыг шаардана. Түгжрэлээс ялгаатай нь хомсдолд мөчлөг (cycle) гэх мэт бүтцийн шинж тэмдэг байхгүй.

\textbf{Практикт:} хүлээлтийн хугацааг хянаж, тогтоосон хязгаараас хэтэрвэл "хомсдолд орсон байж магадгүй" гэж үзэж болно (heuristic). Хомсдлыг шийдвэрлэх арга:
\begin{itemize}[leftmargin=1.5em]
  \item \textbf{Хөгшрөлт (aging):} удаан хүлээсэн урсгалын эрэмбийг аажмаар нэмэгдүүлнэ (5-р бүлэг).
  \item \textbf{FIFO / ээлжийн дараалал:} нөөцийг хүсэлт ирсэн дарааллаар өгөх — хязгаарлагдмал хүлээлт (bounded waiting).
  \item \textbf{Түгжрэлээс сэргэхэд:} victim сонгохдоо урсгалыг хэдэн удаа буцаасныг (rollback) зардлын хүчин зүйлд оруулж, нэг урсгалыг дахин дахин сонгохгүй байх (§8.8.2).
  \item Exercise 8.31 / \texttt{bridge.c fair} горим: нэг талаас дараалан хамгийн ихдээ $K$ хүн оруулаад ээлжийг солих.
\end{itemize}

%---------------------------------------------------------------
\question{8.8}{Бодлого: нөөцийн хүсэлтийг шууд хангаж чадахгүй бол хүлээж (blocked) буй урсгалуудаас хүссэн нөөцийг нь булааж хүсэгч урсгалд өгнө; хүлээж буй урсгалын хүлээх вектор булаагдсан нөөцөөр нэмэгдэнэ. (a) Түгжрэл үүсч болох уу? (b) Хязгааргүй хаагдал (indefinite blocking) үүсч болох уу?}

\textbf{(a) Түгжрэл үүсэхгүй.} Энэ бодлого \textbf{no preemption} нөхцөлийг зөрчдөг: хаагдсан (blocked) урсгалын эзэмшиж буй нөөцийг булааж авдаг. Хаагдсан урсгал нөөцөө "барьсаар" хүлээж чадахгүй тул тойрог хүлээлт тогтвортой хэвээр үлдэхгүй — хэрэв $T_i$ нь $T_j$-ийн нөөцийг хүлээж, $T_j$ хаагдсан бол $T_i$ тэр нөөцийг шууд булааж авна. Ажиллаж буй (хаагдаагүй) урсгалын нөөцийг булаахгүй ч ажиллаж буй урсгал хүлээхгүй тул мөчлөгт оролцохгүй.

\textbf{(b) Тийм, хязгааргүй хаагдал (хомсдол) үүсч болно.} Хаагдсан урсгалаас нөөц байнга булаагдаж болно: жишээнд $T_0$ $(0,0,1)$-ийг хүлээж хаагдсан, тэр хооронд $T_2$ түүний нөөцийг авсан. Хэрэв бусад урсгалууд үргэлжлэн шинэ хүсэлт гаргаж $T_0$-ийн нөөцийг дахин дахин булааж байвал $T_0$-ийн хэрэгцээ хэзээ ч бүрэн хангагдахгүй — хүлээх вектор нь улам томорсоор байна.

%---------------------------------------------------------------
\question{8.9}{Системийн агшин зураг:
\begin{center}
\begin{tabular}{lcc}
\toprule
 & Allocation (A B C D) & Max (A B C D)\\
\midrule
$T_0$ & 3 0 1 4 & 5 1 1 7\\
$T_1$ & 2 2 1 0 & 3 2 1 1\\
$T_2$ & 3 1 2 1 & 3 3 2 1\\
$T_3$ & 0 5 1 0 & 4 6 1 2\\
$T_4$ & 4 2 1 2 & 6 3 2 5\\
\bottomrule
\end{tabular}
\end{center}
Банкирын алгоритмаар дараах төлөв бүр аюулгүй эсэхийг тогтоо: (a) Available $=(0,3,0,1)$; (b) Available $=(1,0,0,2)$.}

Need $=$ Max $-$ Allocation: $T_0 = (2,1,0,3)$, $T_1 = (1,0,0,1)$, $T_2 = (0,2,0,0)$, $T_3 = (4,1,0,2)$, $T_4 = (2,1,1,3)$.

\textbf{(a) Available $=(0,3,0,1)$ — аюулгүй БИШ.}
\begin{center}
\begin{tabular}{lll}
\toprule
Урсгал & Need $\le$ Work? & Work\\
\midrule
$T_2$ & $(0,2,0,0) \le (0,3,0,1)$ $\checkmark$ & $(3,4,2,2)$\\
$T_1$ & $(1,0,0,1) \le (3,4,2,2)$ $\checkmark$ & $(5,6,3,2)$\\
$T_3$ & $(4,1,0,2) \le (5,6,3,2)$ $\checkmark$ & $(5,11,4,2)$\\
$T_0$ & $(2,1,0,3)$: D $3 > 2$ $\times$ & —\\
$T_4$ & $(2,1,1,3)$: D $3 > 2$ $\times$ & —\\
\bottomrule
\end{tabular}
\end{center}
$T_2, T_1, T_3$ дууссаны дараа D нөөц ердөө 2 үлдэх ба $T_0$, $T_4$ тус бүр 3 D шаарддаг тул аль нь ч дуусч чадахгүй. Иймд бүх урсгал дуусах дараалал байхгүй — төлөв аюулгүй биш.

\textbf{(b) Available $=(1,0,0,2)$ — аюулгүй.}
$T_1$: $(1,0,0,1) \le (1,0,0,2)$ → Work $=(3,2,1,2)$;
$T_2$: $(0,2,0,0) \le$ → $(6,3,3,3)$;
$T_3$: $(4,1,0,2) \le$ → $(6,8,4,3)$;
$T_4$: $(2,1,1,3) \le$ → $(10,10,5,5)$;
$T_0$: $(2,1,0,3) \le$ → $(13,10,6,9)$.
Аюулгүй дараалал: $\langle T_1, T_2, T_3, T_4, T_0\rangle$. (\texttt{banker.py}-ээр шалгасан.)

%---------------------------------------------------------------
\question{8.10}{Аюулгүй байдлын алгоритмын кодоо ашиглаж, $\text{Max}_i = \text{Waiting}_i + \text{Allocation}_i$ гэж тодорхойлж түгжрэл илрүүлэх алгоритм хэрэгжүүлж болох уу?}

\textbf{Тийм.} Түгжрэл илрүүлэх алгоритм (§8.7.2) нь аюулгүй байдлын алгоритмтай бараг ижил бүтэцтэй бөгөөд \textit{Need}-ийн оронд одоогийн \textit{Request} (хүлээж буй) матрицыг ашигладаг. $\text{Max}_i = \text{Waiting}_i + \text{Allocation}_i$ гэж тодорхойлбол $\text{Need}_i = \text{Max}_i - \text{Allocation}_i = \text{Waiting}_i$ болох тул аюулгүй байдлын алгоритм яг "одоо хүлээж буй хүсэлтээ авбал урсгал дуусч чадах уу?" гэдгийг шалгана — энэ нь илрүүлэх алгоритмын агуулга.

Ялгааг анхаарах: илрүүлэх алгоритмд $\text{Allocation}_i = 0$ урсгалыг эхнээсээ \textit{Finish}$=$true гэж үздэг (энэ нь $\text{Waiting}_i \le$ Work нөхцөлтэй ихэвчлэн ижил үр дүн өгнө). Мөн тайлбар нь өөр: аюулгүй байдлын алгоритм "үгүй" гэвэл \emph{аюулгүй биш төлөв} (түгжрэл үүсч \emph{болзошгүй}), харин энд "үгүй" бол \textit{Finish}$=$false урсгалууд \emph{одоо} түгжрэлд орсон гэсэн үг — учир нь урсгалууд хүлээж буй хүсэлтээсээ илүүг шаардахгүй гэсэн өөдрөг таамаглалтай.

%---------------------------------------------------------------
\question{8.11}{Зөвхөн нэг урсгалтай (single-threaded) нэг процессоос бүрдсэн түгжрэл байж болох уу? Тайлбарла.}

\textbf{Тодорхойлолтын дагуу үгүй.} Түгжрэл үүсэхэд \emph{circular wait} буюу хоёр ба түүнээс дээш урсгал бие биенийхээ эзэмшсэн нөөцийг хүлээх хэрэгтэй. Нэг урсгалтай процесс нөөцийг эзэмшиж байх үедээ хүлээж байвал, тэр нөөцийг чөлөөлөх өөр урсгал байхгүй тул энэ нь мөчлөг биш.

\textbf{Гэхдээ практикт өөртөө гацах (self-deadlock) боломжтой:} жишээ нь рекурсив биш мутексийг аль хэдийн эзэмшсэн урсгал дахин \texttt{pthread\_mutex\_lock()} дуудвал үүрд хүлээнэ; эсвэл семафорыг 0 үед \texttt{wait()} хийгээд \texttt{signal()}-ийг өөрөө дараа нь хийхээр бичвэл. Энэ нь программын алдаа бөгөөд нөөц хуваарилалтын графын утгаар хоёр урсгалын түгжрэл биш.

%---------------------------------------------------------------
\section{Programming Problems (бичлэгт)}

\textbf{8.32 — нэг эгнээтэй гүүр (\texttt{lab7/bridge.c}).} POSIX урсгал: хойд ба өмнөд зүгийн тариаланч тус бүр урсгал. \texttt{pthread\_mutex} ба \texttt{pthread\_cond}: гүүрэн дээр зөвхөн нэг чиглэлийнхэн байна; эсрэг чиглэлийнх гүүр хоосрохыг хүлээнэ — эсрэг талын хоёр тариаланч зэрэг гүүрэнд орох (түгжрэл) боломжгүй. Гүүрээр явах хугацаа санамсаргүй (\texttt{usleep}). \texttt{fair} горим (Exercise 8.31): нэг талаас дараалан хамгийн ихдээ 2 хүн орсны дараа эсрэг талд ээлж өгч хомсдлоос сэргийлнэ. Программ төгсгөлд "эсрэг чиглэл зэрэг гүүрэн дээр: 0 удаа" гэж шалгана; 30 удаагийн давталтад 0.

\textbf{8.33 — түгжээг дарааллаар авах (\texttt{lab7/Transaction.java}).} Figure 8.7-ийн \texttt{transaction(from, to)}: хоёр урсгал $A\to B$ ба $B\to A$ зэрэг хийвэл тойрог хүлээлт үүсч гацна (\texttt{naive} горим 3 с-ийн дараа "DEADLOCK" гэж илрүүлдэг). Засвар: \texttt{System.identityHashCode()}-ийн бага утгатай дансны түгжээг эхэлж авна; хэш тэнцүү бол нэмэлт \texttt{tieLock} авна. Үр дүн: 400 гүйлгээ бүгд дуусч, нийт үлдэгдэл 2000 хэвээр.

\end{document}
